\chapter{2007年上海卷（秋理）}
\section{填空题}
\begin{problem}
函数 \( f(x)=\frac{\lg(4-x)}{x-3} \) 的定义域是 \fillinblank{} \fillinblank{}.
\end{problem}

\begin{problem}
若直线 \(l_{1}: 2x + my + 1 = 0\) 与直线 \(l_{2}: y = 3x - 1\) 平行，则 \(m = \)  \fillinblank{}.
\end{problem}

\begin{problem}
函数 \( f(x)=\frac{x}{x-1} \) 的反函数 \( f^{-1}(x)= \)  \fillinblank{}.
\end{problem}

\begin{problem}
方程 \( 9^x - 6 \cdot 3^x - 7 = 0 \) 的解是 \fillinblank{} \fillinblank{}.
\end{problem}

\begin{problem}
若 \(x, y \in \R^+\)，且 \(x + 4y = 1\)，则 \(x \cdot y\) 的最大值是 \fillinblank{} \fillinblank{}.
\end{problem}

\begin{problem}
函数 \( y = \sin\left(x + \frac{\pi}{3}\right)\sin\left(x + \frac{\pi}{2}\right) \) 的最小正周期 \(T =\) \fillinblank{} \fillinblank{}.
\end{problem}

\begin{problem}
在五个数 \(1, 2, 3, 4, 5\) 中，若随机取出三个数字，则剩下两个数字都是奇数的概率是 \fillinblank{}.（结果用数值表示）
\end{problem}

\begin{problem}
以双曲线 \(\frac{x^2}{4} - \frac{y^2}{5} = 1\) 的中心为焦点，且以该双曲线的左焦点为顶点的抛物线方程是 \fillinblank{} \fillinblank{}.
\end{problem}

\begin{problem}
对于非零实数 \(a, b\)，以下四个命题都成立：

\begin{circlelist}
\item \(a + \frac{1}{a} \neq 0\)；
\item \((a + b)^2 = a^2 + 2ab + b^2\)；
\item 若 \(|a| = |b|\)，则 \(a = \pm b\)；
\item 若 \(a^2 = ab\)，则 \(a = b\).
\end{circlelist}

那么，对于非零复数 \(a, b\)，仍然成立命题的所有序号是 \fillinblank{} \fillinblank{}.
\end{problem}

\begin{problem}
在平面上，两条直线的位置关系有相交、平行、重合三种. 已知 \(\alpha, \beta\) 是两个相交平面，空间两条直线 \(l_{1}, l_{2}\) 在 \(\alpha\) 上的射影是直线 \(s_{1}, s_{2}\)，\(l_{1}, l_{2}\) 在 \(\beta\) 上的射影是直线 \(t_{1}, t_{2}\). 用 \(s_{1}\) 与 \(s_{2}, t_{1}\) 与 \(t_{2}\) 的位置关系，写出一个总能确定 \(l_{1}\) 与 \(l_{2}\) 是异面直线的充分条件：\fillinblank{}.
\end{problem}

\begin{problem}
已知 \(P\) 为圆 \(x^{2} + (y - 1)^{2} = 1\) 上任意一点（原点 \(O\) 除外），直线 \(OP\) 的倾斜角为 \(\theta\) 弧度，记 \(d = |OP|\). 在右侧的坐标系中，画出以 \((\theta, d)\) 为坐标的点的轨迹的大致图形为

\bitmapfigure[max width=0.75\linewidth,max totalheight=6.0cm]{img/2007/shanghai_science/q11_fig1.png}
\end{problem}

\section{单选题}
\begin{problem}
已知 \(a, b \in \R\)，且 \(2 + ai, b + i\)（\(i\) 是虚数单位）是实系数一元二次方程 \(x^{2} + px + q = 0\) 的两个根，那么 \(p, q\) 的值分别是（\quad）

\choices
    {\(p = -4, q = 5\)}
    {\(p = -4, q = 3\)}
    {\(p = 4, q = 5\)}
    {\(p = 4, q = 3\)}
\end{problem}

\begin{problem}
设 \(a, b\) 是非零实数，若 \(a < b\)，则下列不等式成立的是（\quad）

\choices
    {\(a^2 < b^2\)}
    {\(ab^{2} < a^{2}b\)}
    {\(\frac{1}{ab^2} < \frac{1}{a^2b}\)}
    {\(\frac{b}{a} < \frac{a}{b}\)}
\end{problem}

\begin{problem}
直角坐标系 \(xOy\) 中，\(\bs{i}, \bs{j}\) 分别是与 \(x, y\) 轴正方向同向的单位向量. 在直角三角形 \(ABC\) 中，若 \(\overrightarrow{AB} = 2\bs{i} + \bs{j}\)，\(\overrightarrow{AC} = 3\bs{i} + k\bs{j}\)，则 \(k\) 的可能个数是（\quad）

\choices
    {1}
    {2}
    {3}
    {4}
\end{problem}

\begin{problem}
设 \(f(x)\) 是定义在正整数集上的函数，且 \(f(x)\) 满足：“当 \(f(k) \geqslant k^2\) 成立时，总可推出 \(f(k + 1) \geqslant (k + 1)^2\) 成立”. 那么，下列命题总成立的是（\quad）

\choices
    {若 \(f(3) \geqslant 9\) 成立，则当 \(k \geqslant 1\) 时，均有 \(f(k) \geqslant k^2\) 成立}
    {若 \(f(5) \geqslant 25\) 成立，则当 \(k \leqslant 5\) 时，均有 \(f(k) \geqslant k^2\) 成立}
    {若 \(f(7) < 49\) 成立，则当 \(k \geqslant 8\) 时，均有 \(f(k) < k^2\) 成立}
    {若 \(f(4) = 25\) 成立，则当 \(k \geqslant 4\) 时，均有 \(f(k) \geqslant k^2\) 成立}
\end{problem}

\section{解答题}
\begin{problem}
如图，在体积为 1 的直三棱柱 \(ABC - A_{1}B_{1}C_{1}\) 中，\(\angle ACB = 90^{\circ}\)，\(AC = BC = 1\)，求直线 \(A_{1}B\) 与平面 \(BB_{1}C_{1}C\) 所成角.（结果用反三角函数值表示）

\bitmapfigure[max width=0.38\linewidth,max totalheight=4.8cm]{img/2007/shanghai_science/q16_fig1.png}
\end{problem}

\begin{problem}
在 \(\triangle ABC\) 中，\(a, b, c\) 分别是三个内角 \(A, B, C\) 的对边. 若 \(a = 2, C = \frac{\pi}{4}, \cos \frac{B}{2} = \frac{2\sqrt{5}}{5}\)，求 \(\triangle ABC\) 的面积 \(S\).
\end{problem}

\begin{problem}
近年来，太阳能技术运用的步伐日益加快. 2002 年全球太阳电池的年生产量达到 670 兆瓦，年生产量的增长率为 34\%. 以后四年中，年生产量的增长率逐年递增 2\%（如，2003 年的年生产量的增长率为 36\%）.

\begin{enumerate}
\item 求 2006 年全球太阳电池的年生产量（结果精确到 0.1 兆瓦）；
\item 目前太阳电池产业存在的主要问题是市场安装量远小于生产量，2006 年的实际安装量为 1420 兆瓦. 假设以后若干年内太阳电池的年生产量的增长率保持在 42\%，到 2010 年，要使年安装量与年生产量基本持平（即年安装量不少于年生产量的 95\%），这四年中太阳电池的年安装量的平均增长率至少应达到多少（结果精确到 0.1\%）？
\end{enumerate}
\end{problem}

\begin{problem}
已知函数 \( f(x) = x^2 + \frac{a}{x} \ (x \neq 0,\text{常数 } a \in \R) \).

\begin{enumerate}
\item 讨论函数 \( f(x) \) 的奇偶性，并说明理由；
\item 若函数 \( f(x) \) 在 \( x \in [2, +\infty) \) 上为增函数，求 \( a \) 的取值范围.
\end{enumerate}
\end{problem}

\begin{problem}
如果有穷数列 \(a_1, a_2, a_3, \cdots, a_n\)（\(n\) 为正整数）满足条件 \(a_1 = a_n\)，\(a_2 = a_{n-1}, \cdots, a_n = a_1\)，即 \(a_i = a_{n-i+1}\)（\(i = 1, 2, \cdots, n\)），我们称其为“对称数列”. 例如，由组合数组成的数列 \(\mathrm{C}_m^0, \mathrm{C}_m^1, \cdots, \mathrm{C}_m^m\) 就是“对称数列”.

\begin{enumerate}
\item 设 \(\{b_n\}\) 是项数为 7 的“对称数列”，其中 \(b_1, b_2, b_3, b_4\) 是等差数列，且 \(b_1 = 2, b_4 = 11\). 依次写出 \(\{b_n\}\) 的每一项；
\item 设 \(\{c_n\}\) 是项数为 \(2k - 1\)（正整数 \(k > 1\)）的“对称数列”，其中 \(c_k, c_{k+1}, \cdots, c_{2k-1}\) 是首项为 50，公差为 \(-4\) 的等差数列. 记 \(\{c_n\}\) 各项的和为 \(S_{2k-1}\). 当 \(k\) 为何值时，\(S_{2k-1}\) 取得最大值？并求出 \(S_{2k-1}\) 的最大值；
\item 对于确定的正整数 \(m > 1\)，写出所有项数不超过 \(2m\) 的“对称数列”，使得 \(1, 2, 2^2, \cdots, 2^{m-1}\) 依次是该数列中连续的项；当 \(m > 1500\) 时，求其中一个“对称数列”前 2008 项的和 \(S_{2008}\).
\end{enumerate}
\end{problem}

\begin{problem}
我们把由半椭圆 \(\frac{x^2}{a^2} + \frac{y^2}{b^2} = 1\ (x \geqslant 0)\) 与半椭圆 \(\frac{y^2}{b^2} + \frac{x^2}{c^2} = 1\ (x \leqslant 0)\) 合成的曲线称作“果圆”，其中 \(a^2 = b^2 + c^2\)，\(a > 0\)，\(b > c > 0\). 如图，设点 \(F_0, F_1, F_2\) 是相应椭圆的焦点，\(A_1, A_2\) 和 \(B_1, B_2\) 分别是“果圆”与 \(x, y\) 轴的交点.

\begin{enumerate}
\item 若 \(\triangle F_0F_1F_2\) 是边长为 1 的等边三角形，求“果圆”的方程；
\item 当 \(\left|A_1A_2\right| > \left|B_1B_2\right|\) 时，求 \(\frac{b}{a}\) 的取值范围；
\item 连接“果圆”上任意两点的线段称为“果圆”的弦. 试研究：是否存在实数 \(k\)，使斜率为 \(k\) 的“果圆”平行弦的中点轨迹总是落在某个椭圆上？若存在，求出所有可能的 \(k\) 值；若不存在，说明理由.
\end{enumerate}

\bitmapfigure[max width=0.42\linewidth,max totalheight=4.8cm]{img/2007/shanghai_science/q21_fig1.png}
\end{problem}
